\documentclass[12pt,a4paper]{report}

% --- Core Packages ---
\usepackage[utf8]{inputenc}
\usepackage[margin=25mm]{geometry} % CUED required 25mm margins
\usepackage{amsmath, amssymb, amsfonts} % For mathematical typesetting
\usepackage{graphicx} % For importing diagrams
\usepackage{booktabs} % For academic-grade tables
\usepackage{microtype} % Improves justification and spacing
\usepackage{hyperref} % For clickable citations and cross-references
\usepackage{titlesec} % Format chapter headings cleanly

% --- Optional formatting adjustments ---
\hypersetup{
    colorlinks=true,
    linkcolor=black,
    filecolor=black,      
    urlcolor=blue,
    citecolor=black,
}

\begin{document}

% ==============================================================================
% TITLE PAGE
% ==============================================================================
\begin{titlepage}
    \centering
    \vspace*{2cm}
    {\Huge \bfseries Spiking Neural Networks\par}
    \vspace{1cm}
    {\Large Biomimetic 3D Echolocation and Auditory Pathway Dynamics \par}
    \vspace{2cm}
    {\large \textbf{Jack Henry} \par}
    {\large Churchill College \par}
    \vfill
    {\large Department of Engineering \\ University of Cambridge \par}
    \vspace{1cm}
    \vspace{1cm}
    {\large \today \par}
\end{titlepage}

\pagenumbering{roman}

% ==============================================================================
% TECHNICAL ABSTRACT (Max 2 pages, self-contained)
% ==============================================================================
\chapter*{Technical Abstract}
\addcontentsline{toc}{chapter}{Technical Abstract}
\markboth{Technical Abstract}{Technical Abstract}

% CUED Requirement Reminder: Include Name, College, and Project Title here
\noindent
\textbf{Author:} Your Name \\
\textbf{College:} Your College \\
\textbf{Project Title:} Spiking Neural Networks: Biomimetic 3D Echolocation \\
\textbf{Supervisor:} Your Supervisor's Name \\

\noindent
\textit{[Draft this section last. Ensure it details the problem, the bottom-up SNN techniques used, main results of your three parallel pathways, and your core conclusions regarding biological realism as an optimizer.]}

\newpage
\tableofcontents
\newpage

% ==============================================================================
% CHAPTER 0: ACRONYMS & NOMENCLATURE
% ==============================================================================
\chapter*{Acronyms}
\addcontentsline{toc}{chapter}{Acronyms}
\markboth{Acronyms}{Acronyms}

\begin{tabbing}
    \textbf{SNN} \qquad \= Spiking Neural Network \\
    \textbf{LIF} \> Leaky Integrate-and-Fire \\
    \textbf{RF} \> Receptive Field \\
    \textbf{VCN} \> Ventral Cochlear Nucleus \\
    \textbf{IC} \> Inferior Colliculus \\
    \textbf{AC} \> Auditory Cortex \\
    \textbf{CANN} \> Continuous Attractor Neural Network \\
    \textbf{CD} \> Corollary Discharge \\
    \textbf{DSP} \> Digital Signal Processing \\
    \textbf{SNR} \> Signal-to-Noise Ratio \\
    \textbf{ToF} \> Time-of-Flight
\end{tabbing}

\newpage
\pagenumbering{arabic}

% ==============================================================================
% CHAPTER 1: INTRODUCTION & MOTIVATION
% ==============================================================================
\chapter{Introduction \& Motivation}
\label{ch:intro}

\section{The Engineering Problem: 3D Sound Localization}
% Detail the task of locating objects in 3D space using active emission
% Discuss the edge system constraints (size, weight, and power limits of radar/sonar)
\label{sec:problem}

\begin{figure}[htbp]
    \centering
    % \includegraphics[width=0.8\textwidth]{figures/3d_space_localization.pdf}
    \caption{Geometric schematic of 3D target coordinates (Distance, Azimuth, Elevation) defined relative to an active sensor platform.}
    \label{fig:3d_space}
\end{figure}

\section{Conventional Approaches: Classical Radar and Sonar Systems}
% Detail how modern radar/sonar handle spatial detection (matched filters, array signal processing)
% Highlight where these methods struggle on ultra-low power platforms

\section{The Biomimetic Pivot: Echolocation as an Architectural Blueprint}
% Compare bat bio-sonar to artificial sonar
% Discuss power consumption advantages: why sparse spikes save energy
% Framed as: "The biological brain is the existence proof that this can be optimized"

\section{Architectural Philosophy: Bottom-Up Tuning vs. Black-Box Optimization}
% Explain the pivot away from naive black-box SNN training due to diagnostication limits
% Outline your bottom-up, hand-tuned neural dynamics approach 
% Frame how this structured architecture allows training to be used later in a targeted manner

% ==============================================================================
% CHAPTER 2: NEUROMORPHIC PRIMITIVES \& CLASSICAL ALGORITHMIC PARALLELS
% ==============================================================================
\chapter{Neuromorphic Primitives \& Algorithmic Parallels}
\label{ch:primitives}

\section{Neurons and Spiking: The Leaky Integrate-and-Fire Model}
% Define the foundational math of the LIF neuron for the non-specialist assessor
The sub-threshold membrane potential dynamics of the Leaky Integrate-and-Fire (LIF) neuron are governed by:
\begin{equation}
    \tau_m \frac{dV(t)}{dt} = -(V(t) - V_{rest}) + R I(t)
\end{equation}

\section{The Rosetta Stone: Neuronal Structures as Classical Algorithms}
% Explicitly map out your parallels here using a clear table
This section establishes a functional bridge between classical digital signal processing (DSP) abstractions and neuromorphic primitives:
\begin{table}[htbp]
    \centering
    \caption{The Neuromorphic-Algorithmic Rosetta Stone.}
    \begin{tabular}{ll}
        \toprule
        \textbf{Classical DSP Framework} & \textbf{Neuromorphic / SNN Abstraction} \\
        \midrule
        Fast Fourier Transform (FFT) Bank & Receptive Field (RF) Tonotopic Filter Bank \\
        Cross-Correlator / Time-of-Flight & Delay Lines \& Coincidence Detection LIF Bank \\
        Matched Filter Receiver & Tuned Synaptic Weight Distribution \\
        Monte Carlo Estimation / Sampling & Sparse Event-Driven Spike Distributions \\
        Bayesian Brain Probability Densities & Continuous Attractor Population Topography \\
        \bottomrule
    \end{tabular}
    \label{tab:rosetta_stone}
\end{table}

\section{Preliminary Control Experiment: LIF vs. RF Coincidence Detection}
% Present the setup and results of your initial coincidence detection test
% Use these results to justify the selection of the core primitive block for the pipeline

% ==============================================================================
% CHAPTER 3: SYSTEM ARCHITECTURE \& PERIPHERAL FRONTEND
% ==============================================================================
\chapter{System Architecture \& The Peripheral Frontend}
\label{ch:frontend}

\section{End-to-End Pipeline Overview}
% Present your high-level system block diagram tracing the signal flow
\begin{figure}[htbp]
    \centering
    % \includegraphics[width=0.9\textwidth]{figures/system_block_diagram.pdf}
    \caption{Architectural overview of the end-to-end 3D auditory processing pipeline.}
    \label{fig:system_overview}
\end{figure}

\section{The Acoustic Signal Model \& Synthetic Echo Generation}

\section{The Cochlea Model and Time-Gain Compensation}

\section{The Ventral Cochlear Nucleus (VCN) \& The Ssquelch Architecture}
% Address the Amplitude-Latency Trap here
% Explain how the Endbulb of Held abstracts population rate codes to minimize jitter
% Discuss the "Loose Squelch Gate" parameters to avoid Premature Refractory Lockout

% ==============================================================================
% CHAPTER 4: PARALLEL PROCESSING PATHWAYS (THE CORE)
% ==============================================================================
\chapter{Parallel Processing Pathways}
\label{ch:pathways}
% Frame how the front-end stream splits into 3 specialized processing pipelines

\section{Distance Extraction (Temporal Processing)}
    \subsection{Biological Mechanism in the Midbrain}
    \subsection{Neuromorphic Abstraction: The Integrated IC Equation}
    % Insert the unified sub-threshold differential update equation
    \begin{equation}
        V_{IC}[c, k, t] = \beta V_{IC}[c, k, t-1] + w_{sensory} \cdot S_{VCN}[c, t] + w_{facil} \cdot F[c, t] + w_{cd} \cdot S_{CD}[k, t]
    \end{equation}
    % Detail the Rolling Cascade Gate (Additive Asymmetric Facilitation) vs. Multiplicative traps
    % Detail Cross-Correlation via the Corollary Discharge (S_CD)
    \subsection{Isolated Pathway Parameter Tuning \& Validation Results}

\section{Azimuth Localisation (Binaural Disparity Processing)}
    \subsection{Biological Mechanism: Interaural Level \& Time Differences}
    \subsection{Neuromorphic Abstraction and Network Implementation}
    \subsection{Isolated Pathway Parameter Tuning \& Validation Results}

\section{Elevation Estimation (Spectral Shaping)}
    \subsection{Biological Mechanism: HRTF Pinna Filtering \& Spectral Cues}
    \subsection{Neuromorphic Abstraction and Network Implementation}
    \subsection{Isolated Pathway Parameter Tuning \& Validation Results}

% ==============================================================================
% CHAPTER 5: CORTICAL INTEGRATION \& CONTINUOUS ATTRACTOR DYNAMICS
% ==============================================================================
\chapter{Cortical Integration \& Continuous Attractor Dynamics}
\label{ch:cortical}

\section{The Auditory Cortex Layer: Frequency Channel Collapse}
% Explain pooling parallel tonotopic channels into a unified chronotopic axis to maximize SNR

\section{Contrast Enhancement via Spatial Lateral Inhibition}

\section{The Continuous Attractor Neural Network (CANN) Readout}
% Discuss the implementation of CANN for spatial tracking stability and noise filtering
% Frame this as the robust foundational interface for downstream motor-control loops

% ==============================================================================
% CHAPTER 6: SYSTEM-LEVEL PERFORMANCE \& DISCUSSION
% ==============================================================================
\chapter{System-Level Performance \& Discussion}
\label{ch:performance}

\section{End-to-End 3D Spatial Tracking Results}
% Present unified 3D localization accuracy metrics

\section{Robustness Analysis: Stress Tests Under Varying Noise Scales}
% Evaluate system accuracy under aggressive SNRs and echo dropouts

\section{The Architectural Arbitration of the Simmons-Suga Controversy}
% Evaluate your filter-bank framework against the historic academic debate
% Argue how your model achieves sub-microsecond timing precision via population dynamics 
% without requiring raw, unconstrained mathematical phase-coherence.

% ==============================================================================
% CHAPTER 7: CONCLUSIONS \& FUTURE WORK
% ==============================================================================
\chapter{Conclusions \& Future Work}
\label{ch:conclusions}

\section{Summary of Original Contributions}
% Recapitulate major insights:
% 1. Biological Realism as an Engineering Optimizer
% 2. The Formalized Neuromorphic-Algorithmic Rosetta Stone
% 3. Morphology-Driven, Training-Free Intelligent Edge Architecture

\section{Future Work and Next Steps}
% Highlight concepts designed but un-implemented due to time constraints:
% - Evolution to a hybrid, targeted trainable SNN model
% - Scaling network complexity to handle multi-target environments
% - Porting the event-list architecture onto physical Neuromorphic Hardware
% - Closed-loop real-time tracking and flight control loops

% ==============================================================================
% REFERENCES
% ==============================================================================
\begin{thebibliography}{99}
\addcontentsline{toc}{chapter}{Bibliography}

\bibitem{simmons1973}
Simmons, J. A. (1973). \emph{The resolution of target distance by echolocating bats}. Journal of the Acoustical Society of America, 54(1), 157-173.

\bibitem{suga1989}
Suga, N. (1989). \emph{Principles of auditory cortical-mapping in the mustache bat for echolocation}. Annual Review of Neuroscience, 12(1), 131-173.

\end{thebibliography}

% ==============================================================================
% APPENDICES
% ==============================================================================
\appendix
\chapter{Risk Assessment Retrospective}
\label{app:risk}
% CUED Mandatory Requirement: Maximum one side of A4 commenting on Michaelmas Term risk assessment
% Discuss hazards encountered, accuracy of original risk profiling, and future mitigation insights.

\chapter{Algorithmic Implementations \& Pseudo-code}
\label{app:code}
% Place heavy processing listings, event-list array architectures, or pseudo-code blocks here.

\end{document}